% SOURCE_STATEMENT: "What is the global singularity structure inside a black hole formed by gravitational collapse, in particular the relationship between a weakly singular Cauchy horizon, stronger terminal singularities, and the center?"
% Preparation and final self-audit: CONTRIBUTING.md prompts C and D; source/audit artifacts retained in the local physics_statements_2026_10_11 batch.
% !TeX program = lualatex
\documentclass[11pt]{article}
\usepackage[margin=25mm]{geometry}
\usepackage{fontspec}
\setmainfont{Latin Modern Roman}
\usepackage{amsmath,amssymb,mathtools}
\usepackage{unicode-math}
\setmathfont{Latin Modern Math}
\usepackage{xurl,hyperref}
\hypersetup{colorlinks=true,urlcolor=blue}
\setcounter{secnumdepth}{0}
\setlength{\parindent}{0pt}
\setlength{\parskip}{5pt}
\setlength{\emergencystretch}{3em}
\newcommand{\sref}[2]{\hyperlink{source:#1:#2}{[S#2]}}
\newcommand{\cataloguescope}[1]{\paragraph{Recorded target.}#1}
\newcommand{\R}{\mathbb R}
\newcommand{\C}{\mathbb C}
\newcommand{\N}{\mathbb N}
\newcommand{\Z}{\mathbb Z}
\newcommand{\Q}{\mathbb Q}
\newcommand{\F}{\mathbb F}
\newcommand{\eps}{\varepsilon}
\begin{document}
% UnsolvedMath ID 30004809; source catalogue code OWR-8415343-012
\setcounter{section}{30004808}
\section{Interior Singularity Structure of Dynamically Formed Black Holes}\label{30004809}
{\small UnsolvedMath ID 30004809 · Mathematical Physics}
\subsection{Definitions and mathematical statement}

\paragraph{Shared notation.}$\mathbb N=\{1,2,\ldots\}$; the standard number fields have their usual meanings. All additional conventions are stated locally.


A Cauchy horizon is the boundary of the domain determined by initial data. A weakly singular horizon may admit a continuous metric extension while curvature or higher derivatives diverge. Gravitational collapse produces a black-hole interior whose future boundary can include both such null boundaries and stronger terminal singularities.

\paragraph{Catalogue wording retained for traceability.}
What is the global singularity structure inside a black hole formed by gravitational collapse, in particular the relationship between a weakly singular Cauchy horizon, stronger terminal singularities, and the center?

\paragraph{Scope and source qualification.}This is a structural research question. The matter model and regularity notion must be fixed before claiming a theorem for a particular collapse class.
\subsection{Short English statement}
Describe how weak Cauchy-horizon singularities, stronger singularities and the center fit together inside dynamically formed black holes.
\subsection{Sources}
\begin{itemize}
\item[\hypertarget{source:30004809:1}{S1}] ULAM.AI and the UnsolvedMath Contributors, UnsolvedMath, record 30004809 (OWR-8415343-012). Catalogue curation under CC BY 4.0; definitions and English exposition added locally.
\url{https://huggingface.co/datasets/ulamai/UnsolvedMath}.
\item[\hypertarget{source:30004809:2}{S2}] Mathematical Aspects of General Relativity (2021), Oberwolfach Reports, DOI 10.4171/owr/2021/40, pp. 2157--2267. Original report and mathematical context.
\url{https://doi.org/10.4171/owr/2021/40}.
\end{itemize}
\label{end:30004809}
\end{document}
